% (User input window)

\paragraph{Mass splitting and half-order amplification: the camouflageability of finite moment families}
Theorem \ref{thm:pom-finite-integer-moments-cannot-replace-half} states: Even if a finite number of integer-order collision moments are exactly the same, the Hellinger constants at the endpoints of \(\delta\downarrow 0\) may still be different.
The following proposition gives a more structural mechanism: splitting a small mass block will only produce \(O(\varepsilon^2)\) level perturbations on the sum of all integer powers \(p_q\) (\(q\ge 2\)),
But on the half-order quantity \(A_{1/2}(w)=\sum_x\sqrt{w(x)}\), \(\sqrt{M}\) level amplification is produced;
And this perturbation can be absorbed exactly by the \(O(\varepsilon^2)\) fine-tuning of the heavy block via the local reversibility of the power-sum coordinate graph (theorem \ref{thm:pom-power-sum-vandermonde-jacobian}).
From this, the quantitative indiscernibility and the corresponding information lower bound for the sub-exponential moment family are obtained.

\begin{proposition}[Power sum transformation and half-order root amplification under mass splitting]\label{prop:pom-mass-splitting-moment-root-transform}
Let \(X\) be a finite set, \(w\in\Delta(X)\).
Take \(x_0\in X\) such that \(w(x_0)=\varepsilon\in(0,1)\), and take the integer \(M\ge 2\).
Let \(X^{\langle M\rangle}\) be obtained by replacing \(x_0\) with \(M\) new symbols \(\{x_{0,1},\dots,x_{0,M}\}\) and define the splitting distribution \(w^{\langle M\rangle}\in\Delta(X^{\langle M\rangle})\) as
\[
w^{\langle M\rangle}(x)=w(x)\ (x\in X\setminus\{x_0\}),\qquad
w^{\langle M\rangle}(x_{0,j})=\varepsilon/M\ (1\le j\le M).
\]
Remember the power sum \(p_q(w)=\sum_x w(x)^q\) (\(q\ge 1\)) and the half-order quantity \(A_{1/2}(w)=\sum_x\sqrt{w(x)}\).
Then there is a strict identity for any real number \(q\ge 1\)
\[
\boxed{\;
p_q\!\bigl(w^{\langle M\rangle}\bigr)=p_q(w)+\varepsilon^q\bigl(M^{1-q}-1\bigr)\; }.
\]
And the half-order quantity satisfies
\[
\boxed{\;
A_{1/2}\!\bigl(w^{\langle M\rangle}\bigr)=A_{1/2}(w)+\sqrt{\varepsilon}\,(\sqrt{M}-1)\; }.
\]
\end{proposition}
\begin{proof}
The contribution of the split block on \(p_q\) is \(M(\varepsilon/M)^q=\varepsilon^q M^{1-q}\), and the original single-point contribution \(\varepsilon^q\) needs to be deducted to obtain the power sum identity.
The same applies to the half-order quantity: the contribution of the split block is \(M\sqrt{\varepsilon/M}=\sqrt{\varepsilon}\sqrt{M}\), which is obtained by deducting the original contribution \(\sqrt{\varepsilon}\). Certification completed.
\end{proof}

\begin{theorem}[Local camouflage of finite power sums: splitting perturbations can be precisely absorbed by heavy block fine-tuning]\label{thm:pom-moment-camouflage-finite-orders}
Fix the integer \(Q\ge 2\), and take mutually different template vectors \(\bar{\mathbf a}=(\bar a_1,\dots,\bar a_Q)\in(0,1)^Q\) so that \(\sum_{i=1}^{Q}\bar a_i=1\).
For any \(\varepsilon\in(0,1)\), let \(a_i:=(1-\varepsilon)\bar a_i\) and note \(\mathbf a=(a_1,\dots,a_Q)\).
in alphabet
\(
X^{(0)}:=\{1,\dots,Q,0\}
\)
Define the distribution on
\[
w^{(0)}:=\sum_{i=1}^{Q}a_i\,\delta_i+\varepsilon\,\delta_0\in\Delta(X^{(0)}).
\]
For the integer \(M\ge 2\), let
\(
X^{(1)}:=\{1,\dots,Q\}\cup\{0_1,\dots,0_M\}
\)
and define the splitting distribution
\[
\widetilde w^{(1)}:=\sum_{i=1}^{Q}a_i\,\delta_i+\sum_{j=1}^{M}\frac{\varepsilon}{M}\,\delta_{0_j}\in\Delta(X^{(1)}).
\]
Then for any \(q\ge 2\) we have
\[
p_q\!\bigl(\widetilde w^{(1)}\bigr)-p_q\!\bigl(w^{(0)}\bigr)
=\varepsilon^q\bigl(M^{1-q}-1\bigr).
\]
And there exists \(\varepsilon_0>0\) (depending only on \(\bar{\mathbf a}\)) such that: when \(\varepsilon\in(0,\varepsilon_0)\), for any \(M\ge 2\) there is a unique vector \(\mathbf b=(b_1,\dots,b_Q)\) that satisfies
\[
\sum_{i=1}^{Q}b_i=1-\varepsilon,\qquad
\max_{1\le i\le Q}|b_i-a_i|=O(\varepsilon^2),
\]
and order
\[
w^{(1)}:=\sum_{i=1}^{Q}b_i\,\delta_i+\sum_{j=1}^{M}\frac{\varepsilon}{M}\,\delta_{0_j}\in\Delta(X^{(1)})
\]
An exact match is established when
\[
\boxed{\;
p_q\!\bigl(w^{(1)}\bigr)=p_q\!\bigl(w^{(0)}\bigr)\qquad(\forall q=1,2,\dots,Q)\; }.
\]
Furthermore, the difference of half-order quantities satisfies
\[
\boxed{\;
A_{1/2}\!\bigl(w^{(1)}\bigr)-A_{1/2}\!\bigl(w^{(0)}\bigr)
=\sqrt{\varepsilon}\,(\sqrt{M}-1)+O(\varepsilon^2)\; }.
\]
\end{theorem}
\begin{proof}[Proof sketch]
Let \(\widetilde\Phi(\mathbf b)=(p_1(\mathbf b),p_2(\mathbf b),\dots,p_Q(\mathbf b))\), where
\(
p_k(\mathbf b)=\sum_{i=1}^{Q}b_i^k
\)。
According to the theorem \ref{thm:pom-power-sum-vandermonde-jacobian} (taking \(n=Q\)), when \(b_i\) are not equal,
\(
\det \mathrm{D}\widetilde\Phi(\mathbf b)=Q!\prod_{i<j}(b_j-b_i)\neq 0
\), so \(\widetilde\Phi\) is a local diffeomorphism in the neighborhood of \(\mathbf a=(a_1,\dots,a_Q)\).
\par
The matching condition \(p_1(\mathbf b)=1-\varepsilon\) fixes the total mass, while for \(k=2,\dots,Q\) it is given by
\[
p_k(w^{(1)})=p_k(\mathbf b)+\varepsilon^k M^{1-k},\qquad
p_k(w^{(0)})=p_k(\mathbf a)+\varepsilon^k
\]
The required weight target is known to be
\(
p_k(\mathbf b)=p_k(\mathbf a)+\varepsilon^k(1-M^{1-k})
\)。
The right-hand pair \(k\ge 2\) are all \(O(\varepsilon^2)\) level perturbations, so according to the implicit function theorem, there is a unique \(\mathbf b\) that satisfies all constraints in the neighborhood of \(\mathbf a\), and we get \(\max_i|b_i-a_i|=O(\varepsilon^2)\).
\par
The half-order difference is given by
\[
A_{1/2}(w^{(1)})-A_{1/2}(w^{(0)})
=\sum_{i=1}^{Q}\bigl(\sqrt{b_i}-\sqrt{a_i}\bigr)
\ +\ \sqrt{\varepsilon}\,(\sqrt{M}-1)
\]
Decomposition; the first term is given by the \(C^1\) property of \(\sqrt{\cdot}\) on \(0,1)\) with \(|b_i-a_i|=O(\varepsilon^2)\) of \(O(\varepsilon^2)\) magnitude (constant depends on \(\min_i a_i\)). Certification completed.
\end{proof}

\begin{theorem}[Exponential blind spot for sums of subexponential integer powers: half-order quantities can split exponentially under matching]\label{thm:pom-subexp-moments-blindspot-half}
Suppose \(X_m\) is a finite set, \(n_m:=|X_m|\) satisfies \(\log n_m\asymp m\).
Let \(Q(m)\) be any sequence of positive integers satisfying \(\log Q(m)=o(m)\) (sub-exponential moment order upper bound).
Then there are two families of distributions \(w_m,\widetilde w_m\in\Delta(X_m)\) such that it holds for sufficiently large \(m\)
\[
\boxed{\;
p_q(w_m)=p_q(\widetilde w_m)\qquad(\forall q=1,2,\dots,Q(m))\;}
\]
But there is also a constant \(c_0>0\) so that
\[
\boxed{\;
\frac{1}{m}\log\frac{A_{1/2}(\widetilde w_m)}{A_{1/2}(w_m)}\ \ge\ c_0\;}
\]
Thus, half-order quantities separate in the sense of an exponential rate, while the family of integer power sums of any exponential order remains blind to this separation.
\end{theorem}
\begin{proof}[Proof sketch]
Take any constant \(0<\beta<\alpha<\liminf_{m\to\infty}\frac{1}{m}\log n_m\), and let
\[
\varepsilon_m:=e^{-\beta m},\qquad M_m:=\bigl\lfloor e^{\alpha m}\bigr\rfloor,
\qquad Q:=Q(m).
\]
When \(m\) is sufficiently large, there are \(Q+M_m\le n_m\), and \(Q\) heavy coordinates and \(M_m\) light coordinates can be selected in \(X_m\).
Take mutually different weight vectors \(\mathbf a\) (sum \(1-\varepsilon_m\)) on the heavy block, and concentrate the light mass \(\varepsilon_m\) to a single coordinate to get \(w_m\);
Divide the light mass equally into \(M_m\) coordinates, and use the theorem \ref{thm:pom-moment-camouflage-finite-orders} to make \(O(\varepsilon_m^2)\) fine-tuning on the heavy block, so that \(p_q\) exactly matches \(q\le Q(m)\), thus obtaining \(\widetilde w_m\).
\par
By the half-order amplification law of the proposition \ref{prop:pom-mass-splitting-moment-root-transform} combined with the \(O(\varepsilon_m^2)\) error corrected by heavy blocks,
\[
A_{1/2}(\widetilde w_m)-A_{1/2}(w_m)
=\sqrt{\varepsilon_m}\,(\sqrt{M_m}-1)+O(\varepsilon_m^2)
=\exp\!\Bigl(\frac{\alpha-\beta}{2}m+o(m)\Bigr).
\]
Combining this with a rough upper bound on \(A_{1/2}(w_m)\le \sqrt{n_m}\), we get that there exists \(c_0>0\) such that the exponential lower bound shown holds. Certification completed.
\end{proof}

\begin{corollary}[Half-order endpoint constant is not identifiable for the subexponential integer moment family]\label{cor:pom-half-endpoint-constant-subexp-mom-invisible}
Let \(R_w(\delta)\) be the rate-distortion function of diagonal coupling (theorem \ref{thm:pom-diagonal-rate-explicit-param-and-derivative}), and note the small distortion endpoint constant
\(
A_{1/2}(w)^2-1
\)
(Theorem \ref{thm:pom-diagonal-rate-small-distortion-hellinger-endpoint}).
Then under the construction of the theorem \ref{thm:pom-subexp-moments-blindspot-half}, even if the sum of all subexponential integer powers is known
\(\{p_q(w_m):1\le q\le Q(m)\}\)，
The exponential rate of \(A_{1/2}(w_m)^2-1\) still cannot be determined in a consistent sense, let alone the first-order expansion constant of \(\delta\downarrow 0\).
\end{corollary}
\begin{proof}
By the theorem \ref{thm:pom-diagonal-rate-small-distortion-hellinger-endpoint}, the first-order constant of \(\delta\downarrow 0\) is uniquely determined by \(A_{1/2}(w)^2-1\).
And the theorem \ref{thm:pom-subexp-moments-blindspot-half} gives that under the complete matching of the sum of all powers of \(q\le Q(m)\), \(A_{1/2}\) is still separated at the exponential rate.
Therefore any auditing mechanism relying only on finite powers and data cannot consistently determine this endpoint constant under \(m\to\infty\). Certification completed.
\end{proof}

\begin{theorem}[Exponential lower bound of the candidate set: there are still exponential combinatorial degrees of freedom under subexponential moment observation]\label{thm:pom-oracle-bit-lower-bound-subexp-mom}
Under the setting of the theorem \ref{thm:pom-subexp-moments-blindspot-half}, fix a sufficiently large \(m\) and assume that only vectors are observed
\[
\mathcal{M}_m:=\bigl(p_q(w_m)\bigr)_{1\le q\le Q(m)}.
\]
make
\[
\mathcal{W}_m(\mathcal{M}_m):=\Bigl\{w\in\Delta(X_m):\ \bigl(p_q(w)\bigr)_{1\le q\le Q(m)}=\mathcal{M}_m\Bigr\}
\]
is the set of candidate distributions consistent with the observation.
Then there exists a constant \(c_1>0\) such that for sufficiently large \(m\): \(\mathcal{W}_m(\mathcal{M}_m)\) contains at least one pairwise different finite subset whose size satisfies
\[
\boxed{\;
\log_2 N_m\ \ge\ \exp(c_1 m)\;}
\]
Where \(N_m\) is the cardinality of the subset.
Therefore, in the worst case, if one wishes to uniquely recover the half-order endpoint constant or its exponential rate from subexponential moment observations, an exponential amount of additional information (in binary bits) is required.
\end{theorem}
\begin{proof}[Proof sketch]
Take the parameters \(M_m=\lfloor e^{\alpha m}\rfloor\) in the theorem \ref{thm:pom-subexp-moments-blindspot-half}.
Evenly distribute the light mass \(\varepsilon_m\) to any subset \(S\subseteq \((\varepsilon_m,M_m,q)\), independent of the choice of subset \(S\).
Therefore when the heavy block correction is fixed to the same target value on \(q\le Q(m)\) (theorem \ref{thm:pom-moment-camouflage-finite-orders}),
The distributions obtained for different \(S\) are completely consistent on the observation vector \(\mathcal{M}_m\), but the distributions themselves are different.
The number of desirable subsets is at least \(\binom{n_m}{M_m}\), so that
\[
\log_2 N_m\ge \log_2\binom{n_m}{M_m}
\asymp M_m\log_2\frac{n_m}{M_m}
=\exp(\alpha m)\cdot \Theta(m),
\]
Then there exists a constant \(c_1>0\) such that \(\log_2 N_m\ge \exp(c_1 m)\) holds for sufficiently large \(m\). Certification completed.
\end{proof}

\begin{remark}[Endpoint biphase: finite moment closure at zero rate end and strictification barrier at small distortion end]\label{rem:pom-endpoint-biphase-moment-vs-half}
The quadratic take-off curvature of the zero-rate end \(\delta\uparrow\delta_0=1-p_2(w)\) is only determined by \((p_2,p_3)\) (Theorem \ref{thm:pom-diagonal-rate-zero-rate-endpoint-quadratic-curvature}; see Corollary for folding semantics \ref{cor:pom-diagonal-rate-folded-curvature-auditable}).
In contrast, the first-order constant of \(\delta\downarrow 0\) is uniquely determined by \(A_{1/2}(w)^2-1\) (Theorem \ref{thm:pom-diagonal-rate-small-distortion-hellinger-endpoint}), and by Theorem \ref{thm:pom-subexp-moments-blindspot-half} and Theorem \ref{thm:pom-oracle-bit-lower-bound-subexp-mom} gives strict irreplaceability:
Integer power sum observations of any exponential order cannot determine the constant in the sense of an exponential rate. To eliminate the uncertainty, a strictification channel (\(q=\tfrac12\)) or an equivalent amount of exponential external information needs to be introduced.
\par
Under the collaborative semantics of tensor product closure, the multiplication method of endpoint constants (corollary \ref{cor:pom-small-distortion-constant-tensor-multiplicative}) and the uniqueness of its tensor generator (theorem \ref{thm:pom-small-distortion-synergy-unique-additive-generator}) further show that: if the collaborative amplification of small distortion endpoints is required to be uniformly generated by an additive invariant, then the invariant must be Rényi-\(\tfrac12\) entropy \(H_{1/2}\).
Juxtaposed with the kernel version of the RH/GRH square root spectral bound (definition \ref{def:kernel-rh}) on the spectral side, the same index \(\tfrac12\) can be seen as a common anchor for both types of endpoint regimes: the information theoretic endpoint is given the multiplicative generators in \(H_{1/2}\), while the completion side is characterized by the Carath\'eodory positivity of the logarithmic derivative of the horizon to characterize the critical line regime (corollary \ref{cor:xi-single-slice-rh-typing-principle}).
\end{remark}

\begin{proposition}[Auditable endpoint derivative constant term: strictification channel checked out by \texorpdfstring{$R_w'(\delta)$}{R'(delta)}]\label{prop:pom-strictification-missing-auditable-by-rate-derivative}
Let \(R_w(\delta)\) be the rate-distortion function of diagonal coupling (theorem \ref{thm:pom-diagonal-rate-explicit-param-and-derivative}), and note
\(
C_{1/2}(w):=A_{1/2}(w)^2-1
\)。
Then the limit is established when \(\delta\downarrow 0\)
\[
\boxed{\;
\lim_{\delta\downarrow 0}\ \frac{\delta}{\exp\bigl(-R_w'(\delta)\bigr)}=C_{1/2}(w)\;}
\]
Thus the endpoint constant \(C_{1/2}(w)\) (equivalent to the Rényi-\(\tfrac12\) strictification channel) can be directly audited by derivative observations at sufficiently small distortions.
\par
In particular, under the construction of the theorem \ref{thm:pom-subexp-moments-blindspot-half} there are two families of distributions \(w_m,\widetilde w_m\) that satisfy all sub-exponential order integer powers and are completely consistent with the observations, but \(C_{1/2}(w_m)\) and \(C_{1/2}(\widetilde w_m)\) are separated in the exponential rate sense;
Therefore any auditing mechanism relying solely on this type of integer moment data cannot consistently determine the above endpoint derivative constant terms, and the absence of the strictification channel can be detected by estimating the small \(\delta\) behavior of \(\exp(-R_{w_m}'(\delta))/\delta\).
\end{proposition}
\begin{proof}
According to the theorem \ref{thm:pom-diagonal-rate-explicit-param-and-derivative}, the optimal dual parameters satisfy
\(
\exp(-R_w'(\delta))=(1+\kappa(\delta))^{-1}
\)。
And according to the theorem \ref{thm:pom-diagonal-rate-small-distortion-hellinger-endpoint}, when \(\delta\downarrow 0\) we have
\(
\kappa(\delta)\sim C_{1/2}(w)/\delta
\)，
Therefore
\(
\exp(-R_w'(\delta))\sim \delta/C_{1/2}(w)
\), thus obtaining the limit shown.
\par
The last paragraph is obtained by the direct juxtaposition of the theorem \ref{thm:pom-subexp-moments-blindspot-half} and the corollary \ref{cor:pom-half-endpoint-constant-subexp-mom-invisible}: when the integer moment observation is fixed, \(C_{1/2}\) can still be separated, so the endpoint derivative constant term cannot be uniformly determined by this type of observation; conversely, small \(\delta\) The derivative observation of \(C_{1/2}\) can be directly recycled and form an audit certificate. Certification completed.
\end{proof}

\begin{proposition}[Subexponential moment order cannot consistently estimate the half-order exponential rate: exponential growth is a necessary condition]\label{prop:pom-half-exponent-needs-exp-moment-order}
Under the setting of theorem \ref{thm:pom-subexp-moments-blindspot-half}, let
\[
a_{1/2}(w_m):=\frac{1}{m}\log\bigl(A_{1/2}(w_m)^2-1\bigr),
\]
And consider any estimator \(\widehat a_{1/2}(m)\) that depends only on finite moment observations \(\{p_q(w_m):1\le q\le Q(m)\}\).
If it is required to satisfy the consistent error \(|\widehat a_{1/2}(m)-a_{1/2}(w_m)|=o(1)\) in the worst case, then there must be
\[
\boxed{\;\exists\,c>0\ \text{so that}\ \log Q(m)\ge c\,m\ \text{for sufficiently large}\ m\ \text{holds}\;}
\]
That is, the upper bound of the moment order must grow at an exponential rate; any sub-exponential moment order sequence that satisfies \(\log Q(m)=o(m)\) is not enough to guarantee consistent estimation.
\end{proposition}
\begin{proof}
If \(\log Q(m)=o(m)\), then by the theorem \ref{thm:pom-subexp-moments-blindspot-half} there exists \(\widetilde w_m\) such that \(\{p_q\}_{q\le Q(m)}\) is completely consistent but \(A_{1/2}\) is separated at the exponential rate, so \(a_{1/2}\) The limit points of are not unique under the same observation.
Therefore, the error of any estimator that only relies on this observation cannot go to zero in the worst case; therefore, if the consistent error is required to be \(o(1)\), there must be no \(\log Q(m)=o(m)\). Equivalently: there is a constant \(c>0\) such that \(\log Q(m)\ge c\,m\) holds for sufficiently large \(m\). Certification completed.
\end{proof}

\begin{proposition}[Explicit lower bound for polynomial moment estimates: consistent estimation requires \texorpdfstring{$Q(m)$}{Q(m)} to exponential level]\label{prop:pom-half-exponent-explicit-moment-order-lb}
Under the folding semantics, take \(w_m(x)=d_m(x)/2^m\in\Delta(X_m)\), and note \(n_m:=|X_m|\asymp \varphi^{m}\).
make
\[
A_m:=A_{1/2}(w_m)=\sum_{x\in X_m}\sqrt{w_m(x)},
\qquad
a_{1/2}:=\lim_{m\to\infty}\frac{1}{m}\log(A_m^2-1),
\]
And assume \(a_{1/2}>0\).
Consider a class of \textbf{polynomial moment estimates} constructed only from finite power sum data \(\{p_q(w_m)\}_{1\le q\le Q(m)}\):
For each \(m\), take a polynomial \(P_m(t)=\sum_{q=0}^{Q(m)}c_{q,m}t^q\) whose degree does not exceed \(Q(m)\), and let
\[
\widehat A_m:=\sum_{x\in X_m}P_m\!\bigl(w_m(x)\bigr).
\]
If it is required that there is a set of choices such that \(|\widehat A_m-A_m|=o(A_m)\) (so that \(\widehat a_{1/2}(m):=\frac1m\log(\widehat A_m^2-1)\) satisfies \(|\widehat a_{1/2}(m)-a_{1/2}|=o(1)\)),
then there must be
\[
\boxed{\;
\log Q(m)\ \ge\ \bigl(2\log\varphi-a_{1/2}\bigr)m+o(m)\;}
\]
In particular \(Q(m)\) must be exponential.
\end{proposition}
\begin{proof}[Proof sketch]
For any polynomial \(P\) of degree up to \(Q\), there is a constant \(c>0\) by a consistent approximation standard lower bound (Jackson--Bernstein type; Hölder-\(\tfrac12\) regularity for \(f(t)=\sqrt{t}\)) such that
\[
\sup_{t\in[0,1]}\bigl|\sqrt{t}-P(t)\bigr|\ \ge\ \frac{c}{\sqrt{Q+1}}.
\]
Therefore for any \(P_m\) we have
\[
|\widehat A_m-A_m|
\le
\sum_{x\in X_m}\sup_{t\in[0,1]}|\sqrt t-P_m(t)|
\le
n_m\,\sup_{t\in[0,1]}|\sqrt t-P_m(t)|
\]
And derive the optimal possible error from the lower bound to satisfy
\[
\inf_{\deg P_m\le Q(m)}\ |\widehat A_m-A_m|\ \gtrsim\ \frac{n_m}{\sqrt{Q(m)}}.
\]
If \(|\widehat A_m-A_m|=o(A_m)\) is realizable, then there must be \(n_m/\sqrt{Q(m)}=o(A_m)\).
From \(n_m=\exp((\log\varphi+o(1))m)\) and \(A_m=\exp(\frac12 a_{1/2}m+o(m))\) (because \(a_{1/2}>0\)), it can be reduced to
\[
\frac{1}{2m}\log Q(m)\ \ge\ \log\varphi-\frac12 a_{1/2}+o(1),
\]
That is, we get the explicit lower bound shown. Certification completed.
\end{proof}

